\section{Problem formulation and analysis}
\label{sec:analysis}
Let $\mathcal C=\{1,\ldots,N\}$ be a fixed scene carrier and $q$ a query.
The ideal response $r_q^*$ specifies target membership or semantic relevance
over that carrier. An observation supplies a bounded response $y_e$, a geometric
observation operator $A_e$, and query-dependent information about its scope.
A deployed answer is $\widehat a_q=\mathcal R_t(r_q)$, where the output type $t$
determines the readout $\mathcal R_t$. The evidence-to-response problem is to
estimate $r_q$ from the legally available observations without silently adding
negative information or changing the meaning of $\mathcal R_t$.

We distinguish \emph{applicability}, whether an observation refers to the
requested identity, from \emph{spatial authority}, the locations at which it
can provide support or exclusion. Reliability weights the strength of already
admitted evidence. These quantities need not be calibrated probabilities.
Unknown locations supply neither positive nor negative labels; an output
adapter may still need a declared default where no observation is available.

\paragraph{Proposition 1: partial scope does not identify complete extent.}
Suppose the only constraints on binary membership are $r_i^*=1$ for $i\in P$
and $r_i^*=0$ for $i\in N$, with disjoint $P,N$ and a nonempty unknown set
$U=\mathcal C\setminus(P\cup N)$. If no other observation or constraint
distinguishes assignments on $U$, complete membership is not identifiable.

\emph{Proof.} Pick $j\in U$. Two assignments can agree everywhere except at
$j$, taking values zero and one there, while satisfying every supplied
constraint. An estimator given identical evidence returns the same answer
for both and cannot be correct on both. $\square$

This elementary ambiguity result does not assert that the full RGB input is
uninformative. Additional image evidence, native correspondence or learned
membership can resolve an ambiguity. It shows why the complement of a partial
positive observation is not, by itself, an observed negative. Our estimated
completion supplies a hypothesis about that missing scope, not a proof of it.

\paragraph{Proposition 2: false exclusion attenuates direct consensus.}
For one carrier element, let the effective evidence and fixed-prior mass be
$D>0$, numerator $S$, and response $r=S/D\in[0,1]$. Add incorrectly authorized
negative evidence with response zero and total effective mass $u\geq0$,
holding all previous weights, priors and correspondences fixed. Then
\begin{equation}
 r'=\frac{S}{D+u},\qquad r-r'=\frac{ru}{D+u}.
 \label{eq:attenuation}
\end{equation}
If a positive decision uses $r\geq\tau>0$, it flips to negative exactly when
$u>D(r/\tau-1)$.

\emph{Proof.} Zero-valued evidence adds $u$ to the denominator and nothing to
the numerator. Subtraction gives Equation~\ref{eq:attenuation}; rearranging
$S/(D+u)<\tau$ gives the threshold condition. $\square$

The result applies to the direct consensus operator with fixed existing mass,
not to every robust solver or learned member decoder. It holds with correct
geometry: better alignment cannot repair invalid exclusion authority. It also
does not justify deleting valid negatives. Signed-hint controls test that
separate issue empirically.

\paragraph{Counterexample: membership and semantic rendering do not commute.}
Consider a ray with an unselected front Gaussian of opacity $a\in(0,1)$ and
a selected rear Gaussian of opacity $b\in(0,1)$. The normalized full-scene
member response and the opacity of the selected set are respectively
\begin{equation}
 s_{\mathrm{full}}=\frac{(1-a)b}{a+(1-a)b},\qquad
 \alpha_{\mathrm{selected}}=b.
 \label{eq:readout-counterexample}
\end{equation}
With $a=0.9,b=0.2$, they are approximately $0.0217$ and $0.2$.
They produce different answers under the retained semantic threshold $0.5$
and selected-opacity threshold $10/255$. Removing the front Gaussian changes
transmittance; it is not just thresholding a value in the original renderer.
Thus preserving a selection capability requires preserving its objective and
readout, not only its underlying features. This is a constructive counterexample,
not an accuracy guarantee or a claim that all rays disagree.

\paragraph{Consequences for method and evaluation.}
These statements motivate three linked decisions within one mechanism:
admit evidence for the requested identity, estimate its support/exclusion
scope, and form the response with a compatible forward model and readout.
They predict that reconstructing incorrect evidence more accurately can retain
or amplify its errors. They do not predict the sign or magnitude of a benchmark
gain. We test the corresponding design choices with fixed-evidence solver
controls, matched reference-scope interventions and separate semantic/selection
evaluation, rather than treating lower source loss as evidence of correctness.
